\documentclass[11pt]{article}

\usepackage[margin=1in]{geometry}
\usepackage[T1]{fontenc}
\usepackage[utf8]{inputenc}
\usepackage{lmodern}
\usepackage{microtype}
\usepackage{amsmath,amssymb}
\usepackage{xcolor}
\usepackage{enumitem}
\usepackage{titlesec}
\usepackage{xurl}
\usepackage{hyperref}

\DeclareUnicodeCharacter{01C2}{\ensuremath{\ddagger}}
\DeclareUnicodeCharacter{01C1}{\ensuremath{\Vert}}
\DeclareUnicodeCharacter{016B}{\={u}}

\hypersetup{
  colorlinks=true,
  linkcolor=blue!45!black,
  urlcolor=blue!45!black,
  citecolor=blue!45!black,
  pdftitle={Presentation of Hoapa as Pure Form of Communication},
  pdfauthor={Hoapa project}
}

\setcounter{tocdepth}{3}
\setlist[itemize]{leftmargin=1.6em,itemsep=0.35em,topsep=0.35em}
\setlist[enumerate]{leftmargin=1.8em,itemsep=0.35em,topsep=0.35em}
\setlist[description]{leftmargin=1.6em,itemsep=0.55em,topsep=0.35em,style=nextline,font=\normalfont}
\setlength{\parindent}{1.2em}
\setlength{\parskip}{0pt}
\titlespacing*{\section}{0pt}{2.2ex plus 0.6ex minus 0.2ex}{1.0ex}
\titlespacing*{\subsection}{0pt}{1.8ex plus 0.4ex minus 0.2ex}{0.8ex}
\titlespacing*{\subsubsection}{0pt}{1.4ex plus 0.3ex minus 0.2ex}{0.6ex}
\sloppy

\begin{document}
\begin{center}
{\LARGE\bfseries Hoapa protocols as the \emph{pure form} of communication systems\par}
\vspace{0.65em}
{\large Hoapa project, June 2026\par}
\end{center}
\vspace{0.75em}
\tableofcontents
\newpage

\section{Foundations}

The Hoapa project is composed of three layers:

\begin{itemize}

\item a communication system (in the form of a computer program)

\item its verification layer (via a proof-checker)

\item a documentation of what they are in theory and how they develop in practice

\end{itemize}

\subsection{Scope}

The project belongs to a broad category of "informational architectures" which has been studied by disciplines such as:

\begin{itemize}

\item Cybernetics/systems theory

\item Collective intelligence

\item Process Philosophy

\item Protocol studies

\item Semiotics

\item Media ecology

\item Concurrency (computer science)

\item Information theory

\item Non-equilibrium Thermodynamics

\end{itemize}

Our framework assumes the point of view of yet another: \textbf{Applied Category Theory}. Where category theory is the most abstract branch of mathematics, a silent and slow epistemic revolution happening behind closed doors, whose potential to profoundly change the way we perceive reality has been recognized by some insiders, ignored by most, and never really expanded beyond its mathematical origins, despite its universality.

\emph{It's the study of how things relate, regardless of what they are. It asks: "What patterns of connection, transformation, and composition appear everywhere?"}

Applied Category Theory takes those formal patterns and uses them to:

\begin{itemize}

\item Design efficient communication systems

\item Model complex interactions compositionally with guarantees about their behavior

\item Translate between different disciplines, so insights from one domain can inform another

\end{itemize}

Which more or less mirrors the three layers of our project: the system, its verification and its working documentation.

Examples of Hoapa's relatives include: \\ The internet, the universal postal union, supply-chain logistics networks, road signaling systems, distributed computing systems, the genetic code, underground fungal forest networks, flocks of starlings in murmurations, nervous systems, immune systems, Wikipedia, social media, citizen assemblies, the scientific method, conversations\ldots{}

\subsection{Objective}

The main objective of this project is to demonstrate the optimality and accessibility of Hoapa protocols.

We do so by leveraging the universality and expressiveness of category theory to recognize Hoapa as a distinguished element in its domain.

\emph{Category theory is made out of categories, which are like webs of family relations (morphisms) between similar enough things (objects). The existence of a "new" category is often revealed by a bridge (adjunction) between two "prior" categories, and this bridge also identifies certain objects as being the "original incarnation" of the family members (initial object), their "final incarnation" (terminal object), or their "default incarnation" (free object).}

This is the story we're trying to tell here: there is a "pure form" of communication architecture whose advantageous social, engineering or scientific characteristics are reflected in its mathematical properties.

This can be used practically in two ways:

\begin{itemize}

\item The optimality to organize any kind of collective work

\item The "naturality" to start the process uncontroversially

\end{itemize}

Just like many games and sports start at the center of the field of play, Hoapa can be the center from which productive communication starts. (The fact that Hoapa protocols are absolutely decentralized, autonomous and capable of self-modifications to adapt to their context makes them all the more suitable for the role.)

\subsection{Mathematical construction}

The two categories our bridge \textbf{Adj} will connect are:

\textbf{Op:} The operational structure of Hoapa networks \\ \textbf{Ct:} The context in which they are optimal

\subsubsection{Op --- the operational}

The \textbf{Op} side is necessarily an intricate construction because it must navigate several trade-offs to establish decentralization, autonomy, openness to peers, expressiveness, transparency, scalability and self-reflection. But all of them are won with the same strategy: separation of concerns, or modularity---each part of the system must only do one thing.

All the structural properties descend from an original separation which encapsulates the entire system, between:

\begin{itemize}

\item a community of \emph{unityped} communicating agents

\item layers of \emph{multityped} documents produced by the communication of the agents.

\end{itemize}

This has a precise meaning in computer science (known as "untyped/typed"), which is roughly the double specification that:

\begin{itemize}

\item \emph{All participants must "play by the same rules" (like basketball players, and unlike goal keepers vs fields players in sports with goals---or unityped checkers vs multityped chess pieces), meaning the communication system can't distinguish them. To the system, participants are computational resources with all the same characteristics. Whenever it needs resources, it taps into its reserves at random.}

\item \emph{Whatever the participants produce, it should be put in meaningful containers and properly labelled. Those labelled containers describe the data they enclose and guarantee it has certain properties, so we can safely run certain operations on them (if you know you have two numbers, you know you can add them).}

\end{itemize}

Here, contrary to the participants who are all of the same type, the multiplicity of types is the crucial mechanism through which the system copes with \emph{complexity of output}.

Importantly, "unityping" should not be interpreted as "the system wants all participants to be the same". On the contrary, the whole point of this is to welcome and make use of diversity: it is the crucial mechanism through which the system copes with \emph{diversity of input}.

Partial and compact versions of the operational side exist in this world, in the field of participative/direct democracy: citizen assemblies, and they've proved their worth and efficiency time and time again. A small number of participants are selected from the entire electoral register, they're given ample time and resources to study and deliberate, they write a report of their conclusions by consensus only.

Hoapa would be the "networkization" of such units into an open and autonomous protocol. Ideally, it would be embedded in a socio-political context which reflects its philosophy, pacified and circular (unlike the one we have now, but which it can repair by co-evolution, from the ground up).

\subsubsection{Ct --- the contextual}

At the integrated level, the context \textbf{Ct} becomes the recognition that the world is profoundly perspectival: all assertions must be evaluated as points of view, never from a supposedly neutral or absolute standpoint, simply because those don't exist.

Perspectivism is represented mathematically in our framework by \emph{(modal) homotopy type theory}, starting in \emph{sheaf theory}. To its advocates, it's the best foundations we have for the whole of mathematics at the moment, both for theory and practice. It is in particular much more capable of reasoning about relations and equality, while the long established \emph{set theory}, of which it is an extension, focuses on objects viewed in isolation.

\subsubsection{Adj --- the structural relation}

The bridge is then the statement that the generative complexity emerging from such perspectival contexts is what \textbf{allows} concurrent architectures like Hoapa to \textbf{unfold}, while the operational layer of Hoapa is what is \textbf{required} by this context to \textbf{map itself}, and eventually attain recursion. (Where the context is represented by the organization of a community of agents embedded in this context, or how we can view an ecosystem as a network of relations between organisms.)

In a sense, the two ends are quite self-evident:

\begin{itemize}

\item That modalities/perspectivism allow complexity to develop is trivial in both abstract mathematics (Bill Lawvere, Yuri Manin) and philosophy (Gilles Deleuze).

\item That typed data is more robust, readable, literate, composable\ldots{} is trivial in any domain concerned with information transfer, such as programming, whereas the fact that large numbers of unityped agents are most adapted to handle complexity is trivial in concurrency theory.

\end{itemize}

The bridge simply establishes the recognition that those facts are intimately linked.

The relation is still conjectural at this point. But while building for it, we've extracted a first theorem which can be interpreted as "there are no privileged points of view on the operational side of the construction". In mathematical terms it is \emph{a machine-checked adjunction of sheaves over a site of causal interaction contexts, exhibiting every globally coherent state as the free/canonical gluing of its event-local views, and lifting the resource-sensitive logic of the data, in Benton's linear/non-linear sense, up to that same global level from values to sheaves}.

\section{Implementation}

To establish the accessibility and optimality of Hoapa protocols we must also produce and test them at scale.

A certain amount of scientific recognition is necessary to justify the collective efforts needed for an initial implementation---which it can obtain through the clear, fruitful and inspiring theoretical framework this project aspires to be. But the true test of the validity of the conjecture lies in its realization.

\subsection{Stack / programming languages}

While the mathematical verification program is written in Agda---a math-oriented proof-checker designed for homotopy type theory---the usable program which it verifies will be concurrency-oriented and written in:

\begin{itemize}

\item Elixir/Erlang for the part of the program which connects participants---a language based on the same model of concurrency as our framework: the actor model.

\item Rust for the part of the program which orchestrates the transformation of the communication into typed documentation---a low-level "systems language" which is well adapted for the task. It also has the advantage of being used by Pijul, a version control program based on patch theory which is very close in spirit to our protocol, so the two programs can share code.

\end{itemize}

\subsection{Implementation details}

The actor model of concurrency is straightforward, all agents have the same privileges and they can:

\begin{itemize}

\item \emph{send messages}

\item \emph{receive messages and update their internal state accordingly}

\item \emph{create new actors (for instance if they "feel overworked")}

\end{itemize}

Every human participant will be considered an actor in the sense that they can talk/write, listen/read, think and update their worldview accordingly, and tell their friends about the protocol. And the door will always be open for any volunteer to come in, or come out of the protocol.

However, in our framework, the true actors, or "units of computation", are not the individuals but the conversations.

Every time the system needs something to be done it will create one: a small number of registered participants are selected at random and given a specification for the task (in the actor model's view: the protocol sends a request to a newly created conversation, and expects a response from it). And every conversation created by the system can in turn create new conversations to break down its task in smaller increments of work: it can send requests or questions, and expect answers back.

It is in effect an orchestrated conversation of conversations.

\emph{Here the orchestration part is taken care of by session types, another model of concurrency which works in the opposite way as the actor model and brings a new layer to our framework, which tracks and remembers the "shape" of the global conversation, so that the documentation produced by the communication knows its own structure and history: we get traceable data.}

While the orchestration layer can be thought of as the "how" for the data, we use the same strategy for its "what", and add another layer of types (again, labelled containers which describe the data and guarantee certain properties and operations) which in this case gives an executable content-map of what is being discussed.

We repeat the type/container strategy every time we need to add complexity to the system: for its autonomy (typed conversations about the system itself), for its openness and interoperability (typed conversations about other protocols), or for its internal language (how the content types can be given a grammar), etc.

In total there are 6 such layers in the current framework, either encoded or planned. They are, like the unityped agents/typed documentation, all motivated by an engineering trade-off and handled by a mathematical structure:

\begin{itemize}

\item local specificities vs global coherence of content (\emph{sheaf descent})

\item stability vs adaptability of protocol (\emph{protocol-version category \ensuremath{\Pi} with a patch-theoretic fragment category})

\item expressive power vs transparency/accessibility of language (\emph{concept category \ensuremath{\Lambda}, likely 2- or (\ensuremath{\infty},1)-categorical})

\item exploration vs exploitation of concepts (\emph{factorization algebra over the actor site valued in \ensuremath{\Lambda}})

\item internal coherence vs plurality of protocols (\emph{2-category of protocols with Hoapa-respecting geometric morphisms})

\item receiver-relative reading of complexity (\emph{information cohomology})

\end{itemize}

\section{Use cases}

\subsection{A park}

Hoapa can also operate at small scales and with only the "base layer of complexity" which separates the unityped agents from the typed documentation they produce.

For instance we could have a neighborhood "managing" its local park via a protocol in which the motivated residents discuss how it could be enhanced, and produce a to-do list for its maintenance.

Here it is crucial to understand that the system can preserve its advantageous mathematical, social and engineering properties if it is \emph{only} a communication system, and not a full management system which also enacts in reality what has been discussed in the abstract.

Its task is to chart the world, not inhabit and certainly not control it. (That kills all the carefully crafted separations of concerns, starting with the unityping of agents, because it necessarily introduces some hierarchical structure and tension---for basically the same reasons you can share an idea with the entire world without losing it, but you can't do that with a loaf of bread.)

Hoapa protocols can interact with the world only behind buffers which respect individual freedom:

\begin{itemize}

\item If it needs something it simply asks for it, and expects those needs to be met (for instance, if it uses the internet, it states via a self-reflective conversation how much bandwidth it needs, and specifies the characteristics of its network). Individuals or communities, outside of the protocol, would then respond to those needs without compensation.

\item If it wants something to be known it simply makes it as available and accessible as possible, which includes a pedagogical effort in presentation and terminology. But it can't advertise, at the risk of favoring certain targets over others.

\item In some cases, if it wants something to be done it also asks for it, by creating a conversation of a special type, whose members, still selected at random, are charged with the task of pushing for the realization of the desired outcome, again without compensation for themselves or the specialists they convince to do the work. (But philosophically, protocols "prefer" to focus on the documentation and let free individuals and communities decide what they want to do with it.)

\end{itemize}

In either direction, we can think of Hoapa as a coach in sports: its domain of action is preparation and strategy, but it can never enter the field of play or control the players. Only the players play.

In the park example, we can imagine an institution/Hoapa hybrid in which the city government has delegated the supervision of the park to the protocol, but it's still the paid workforce of gardeners which will perform the tasks suggested by the protocol. But ideally, in the spirit of the protocol, it would be only voluntary work, suggested by the protocol but decided at the level of the individuals and communities.

Which makes Hoapa protocols especially well suited to organize the open source economy.

\subsection{Open source}

Because open source produces code, it "lives in the abstract" entirely. Historically, and by definition, only it has only used voluntary labour.

Hoapa is especially adapted to resolve the main recurring problem open source projects face: scalability and continuation beyond the circle of core contributors, usually the instigators of the project. Because every open source project is based on this "circles of knowledge/influence" architecture, which is in essence a hierarchy, it is destined to terminate the moment the inner circle quits (which, as humans, they're bound to do eventually).

Hoapa revokes the privileges of the inner circle, distributes the workload evenly and transfers all the essential knowledge from the original designer's mind into a strictly typed and operational documentation. (As in \emph{literate programming}, by Donald Knuth, the great missed opportunity of software development.)

While this is more demanding in the early stages of the process, it pays off by making the project truly scalable and persistent (not persistent in the sense of a capitalistic endeavour, which always needs to grow and colonize new markets, but closer to how ecosystems work, with minimalism and circularity).

The same arguments can be used to suggest a reform of how science is conducted.

\subsection{Science}

Universities are extremely hierarchical political machines whose ties to reality are clearly not behind buffers, which impair their capacity to describe the world they inhabit.

While the scientific method prevents them from stating things that are wrong in the domains it is applicable, it doesn't prevent a mismanagement of efforts towards certain technologies at the expense of understanding, the fragmentation of the study of Nature into opposing departments, the proliferation of opaque jargons, etc.

Hoapa protocols, on the other hand, are ideal for the realization of what citizen science promised to be, but only marginally developed. In a very real and demonstrable sense, they are the extension of the Galilean scientific method.

\emph{This method defines how one scientist is supposed to ask questions to Nature if they want to receive a reliable answer: through rigorous empirical observation and controlled experimentation, translating those observations into precise mathematical models to test hypotheses.}

In our protocol that would be the types of the orchestration layer, here transposed to a relation to Nature.

In the Hoapa extension, we also define how scientists are supposed to converse among themselves, how they should refine information and eventually store it into a strictly typed and operational technical language which, like the pan-genome is the emanation of chemistry at the surface the planet, would be the reflection of the world we inhabit in language. (Importantly, we can use the amazing properties of homotopy theory to federate all jargons into a "learn once, use everywhere" terminology.)

Practically, one has to visualize a system in which the randomly selected participants, when faced with a technical problem, find and contact the specialists of the domain in question. This is exactly what happens in all citizen assemblies, and again, it's working extremely well in that setting, with both parties reporting high levels of engagement, communication and commitment to understanding.

However, when specialists imagine a world in which the main mode of organization doesn't recognize specialization, they conclude that all the benefits of specialization are lost. But it's not recognition that made them experts in the first place, it's education. And Hoapa protocols are designed for education, as part of their mandate to treat everybody on an equal footing. What is removed with the status is not the knowledge, but the hunt for the status, and at the same time it resituates the obligation to make one's specialized knowledge accessible to others.

Crucially, each individual practitioner gets to decide what questions are worth pursuing.

\subsection{Mathematics}

All of the above applies in mathematics as well, but here things get even more interesting, because the protocol itself is built on the very concepts that its practitioners claim is the correct foundations for the whole field (the aforementioned homotopy type theory and its modal variant, as the internal language of the more general setting of an \ensuremath{\infty}-topos).

This could accomplish one of the many dreams of the main architect behind those structures, Alexander Grothendieck, whose memoirs weave together reflections on creativity, mathematical practice and mathematics itself.

The proof assistant we use in our framework, Agda, is a direct descendant of this history.

\subsection{Generalized permaculture}

I believe Hoapa protocols can stage and target the kind of recursive and self-reflective patterns just mentioned, in the ecosystem, by generalizing the concept of permaculture and applying it to a vast majority of human endeavors. Permaculture could in effect become the default mode of being of humanity, following this very simple algorithm:

\begin{itemize}

\item Make a very detailed and executable map of your environment, including you (where Hoapa protocols shine).

\item Use it to figure out what is best for both you and your environment, and do it (which becomes the default if individual freedom replaces "investability" as the condition for action).

\end{itemize}

This is where the "layers of complexity" which guarantee perspectivism, decentralization, autonomy, openness to peers, expressiveness, transparency, scalability and self-reflection, start to be really useful.

\subsection{Care}

Care can be defined rigorously by another simple 2-step algorithm, which can be viewed as a strategy for networks to ensure their continuity in time:

\begin{enumerate}

\item \emph{In a given network, the evolution towards maturity of newer elements is taken care of by the mature elements already present in the network}

\item \emph{In turn, those newer elements become mature and take care of a new generation of elements}

\end{enumerate}

\emph{Which is an abstraction of John Bowlby's theory of attachment, a generalization at the interdisciplinary level of a process happening in the domains of biology, psychology and evolution.}

Because it is a recursive and self-reflective pattern (step 1 = step 2, or more appropriately 1 \ensuremath{\leftrightarrow} 2), we also find traces of \emph{caring} in the technicalities of our framework:

\begin{itemize}

\item on the side of perspectivism, where we model the context in which the agents operate, in what makes this structure cyclical (\emph{the unit of the adjunction})

\item on the operational side, where we model the conversation between agents, in what makes this structure reciprocal (\emph{operadic boundary-swap})

\end{itemize}

Hoapa can serve as the "\ensuremath{\leftrightarrow}" in "1 \ensuremath{\leftrightarrow} 2", and in fact our conjecture is that it is both required and optimal for this caring algorithm.

\subsection{Networks}

In all the preceding examples we've looked at instances of Hoapa protocols in isolation, but each of them would be better envisioned as nodes of a vast network.

A mature protocol regulating a park would necessarily have typed conversations with other park protocols, especially those close by, to share information and reuse each other's research. It would also have ties to all the independent but interlinked networks a city would be made out of: electricity grid, water management, transportation, agriculture..., and then beyond the city limits, into all sorts of possible protocols-as-points-of-view entities, from playground area design to natural philosophy.

It would also use the code of several open source projects to manage its documentation, and information coming from scientific protocols to apply the precepts of generalized permaculture on its grounds.

In all of this, it would be both an agent and a recipient of care, as a concrete ecosystem as well as an abstract community of participants.

\subsection{Justice and police}

We've listed some of the most positive aspects of life in this presentation, but Hoapa, because it has all these buffers with reality built-in, is also well adapted to coordinate efforts in the regulation of difficult situations, without compromising freedom.

Again we can represent the Hoapa philosophy with a two step algorithm:

\begin{enumerate}

\item The individual chooses to opt into one of the multiple justice-protocol available

\item Each justice-protocol signs up for a police-protocol

\end{enumerate}

All individuals select the authority to which they want to respond, which exists as a documentation and can only be "instantiated" by a jury of peers (a random selection of people, just like in many existing judicial systems; in Hoapa, a typed conversation). The common legal system of two parties which have signed up for different justice protocols is a "\emph{homotopical}" mixture between the two (preestablished by the typed conversations of the "openness to others" layer of each protocol).

This and other self-regulating properties are guaranteed by a combination of attributes of Hoapa protocols, again starting with the unityped-agents/typed-documentation separation of concern: perspectivism, decentralization, autonomy, multiplicity and openness to peers\ldots{}

(Those ideas were put forward by Emmanuel Dockès, a professor of law whose 'Voyage en Misarchie' has been a huge influence on my unlearning path to accept freedom.)

\section{Related structures}

\subsection{Experiments}

While we can find in the history of all the domains just mentioned many examples of projects targeting this state of affairs "philosophically", they never did so with the unityped/typed discipline of Hoapa and its verification layer, which is essential to not see the "old habits" of hierarchical structures reemerge.

This is the fate of most attempts at community-driven structures, especially when they were most active in the post WWII era until the 1970s. They would only exchange the formal hierarchical structures for informal ones, because the root cause of their emergence wasn't properly diagnosed and addressed at the organizational level.

In Hoapa, communities do the "dirty work" in a special place under strict conditions, voluntarily, in order to organize networks of solidarity, so that eventually all individuals can have both freedom and dignity in their lives.

As for the argument that hierarchical structures are part of "human nature", it simply doesn't match our history and evolution, contrary to deep-seated assumptions which only look at a partial rendering of the tail end of the story and often dress Nature in a grotesque 'only the strong survive' frame, which is the antithesis of its autonomous and distributed operational core---a worldview best represented by Herbert Spencer and still extremely influential today, especially in biology and economics.

The lion has never been a king. Humans, on the other hand, have a \textasciitilde{}300,000 years long history of autonomous and distributed cooperative behavior, whose most outstanding payoff is certainly natural language. Hoapa suggests we keep on working on it, in those working conditions.

\subsection{Theories}

\begin{itemize}

\item Charles S. Peirce's --- community of inquiry

\item Gordon Pask --- conversation theory

\item Deleuze / Guattari --- Rhizome \& plane of immanence

\item Maturana / Varela --- autopoiesis

\item Robert Rosen --- (M/R) systems (metabolic/repair)

\item Barr / Chu - Chu spaces

\item Karl Friston --- active inference \& free-energy principle

\item Alain Connes --- noncommutative geometry \& modular flow

\item Bennequin / Belfiore / Laforgue (Huawei) --- topos theory for cognition

\end{itemize}

\section{Name}

'Hoapa' means "everywhere" (every domain) in Khoekhoegowab (of the Khoe language family, spoken in Namibia, Botswana, and South Africa primarily by three ethnic groups: Namakhoen, ǂNūkhoen, and Haiǁomkhoen), to be pronounced like in English, and chosen for 3 reasons:

\begin{itemize}

\item The surface area of the concept to be named is so large, modular and decentralized that it's not possible to select an element or a viewpoint and call it the essence of the system. The alternative is to point to the very multiplicity of the concept for its name.

\item As an homage to San culture, which has been a wonderful source of inspiration for this work. In many ways they are the last known human practitioners of Hoapa, along with the Hadza and other hunter-gatherers of the Congo Basin, so it's only natural the name should come out of one of their languages.

\item It is easy and new to the general public and research community alike, so no one will be jumping to conclusions when they encounter it.

\end{itemize}

\section{Summary of Hoapa protocols in four steps}

\subsection{Eternal}

The protocol has certain \textbf{invariant} mathematical properties which translate into engineering features as a system, into social dynamics as a communication system, and even allows students of Nature to recognize its signature in many systems occurring at all scales in the universe.

\subsection{Initial}

Those properties make it particularly adapted to be the \textbf{starting point} of any collective endeavour, because it has this "default configuration" aspect that favors no one, and this "initial configuration" which is the most welcoming to future change and complexity.

\subsection{(Net)work}

The protocol is then the place where communication \textbf{unfolds} to collect and refine information about a shared context, and produce a useful output in the form of a documentation. Both the protocol and the documentation think of themselves as being part of a network of similar objects, they are constantly working on their external relations as much as their internal growth.

\subsection{Self-evolution}

As it is doing its work of mapping its context, the protocol is also continuously monitoring itself and using its resources to enact its own \textbf{evolution}.

\section{Repository}

\href{https://codeberg.org/doc2cog/hoapa}{The repository} --- Hosted by Codeberg, a truly open source alternative to Microsoft's GitHub.

\subsection{Folder structure}

\begin{description}

\item[\href{https://codeberg.org/doc2cog/hoapa/src/branch/main/formal}{/formal}]
The code files for the verification layer.

\item[\href{https://codeberg.org/doc2cog/hoapa/src/branch/main/communication}{/communication}]
All files related to the presentation of the project (see 7.2).

\item[\href{https://codeberg.org/doc2cog/hoapa/src/branch/main/elixir}{/elixir}]
The code files for the protocol's module concerned with communication between agents.

\item[\href{https://codeberg.org/doc2cog/hoapa/src/branch/main/pedagogy}{/pedagogy}]
Mirrors the /formal folder and walks the reader through what the code does chunk by chunk (AI generated and unfortunately largely inaccessible to non-mathematicians).

\item[\href{https://codeberg.org/doc2cog/hoapa/src/branch/main/workflow}{/workflow}]
The files used and written by LLM coding agents to ground and track the progress of the project.

\end{description}

\subsection{Documents of interest}

\begin{description}

\item[\href{https://codeberg.org/doc2cog/hoapa/src/branch/main/communication/linear_nonlinear_adjunction_abstract.pdf}{Extended abstract: "A linear/non-linear adjunction of sheaves\ldots{}"}]
For mathematicians, the first, deliberately modest theorem of the framework, fully machine-checked in cubical Agda.

\item[\href{https://codeberg.org/doc2cog/hoapa/src/branch/main/communication/adjunction_diagram.png}{Diagram: Construction map for the cross-side adjunction}]
A diagram presenting to mathematicians the conjectured adjunction, the bridge this project is trying to establish.

\item[\href{https://codeberg.org/doc2cog/hoapa/src/branch/main/communication/reference_zalamea.pdf}{Letter of Recommendation: From Fernando Zalamea}]
From Fernando Zalamea, a mathematician and philosopher, leading expert on the work of Alexander Grothendieck and Charles S. Peirce.

\item[\href{https://codeberg.org/doc2cog/hoapa/src/branch/main/communication/general_evolution.pdf}{Article: "An Introduction to Grothendieck’s General Theory of Evolution"}]
For exact and natural scientists, about the dynamics our universe keeps on repeating, a sort of Hoapa-informed version of Philip Anderson's "More is Different".

\item[\href{https://codeberg.org/doc2cog/hoapa/src/branch/main/communication/univocal_science.pdf}{Article: "Univocal Science --- From Spinoza to Grothendieck, missed opportunities for a category-theory-driven scientific revolution"}]
For philosophers and philosophers of mathematics, the "seed" of the conjectured bridge.

\item[\href{https://codeberg.org/doc2cog/hoapa/src/branch/main/communication/liberation_engineering.pdf}{Manifesto: "Liberation engineering for resistance networks"}]
For social scientists, a diagnosis of institutional hierarchy as a "virus" in humanity's world-describing systems, and the argument for decentralized liberation networks as the remedy.

\end{description}

\end{document}
